\clearpage\pagestyle{otMain}\fontsize{11}{15}\selectfont\corrigesChapitre{2}{Puissances, racines et notation scientifique}
\begin{blocLivre}{cor-2-1}\begin{corrige}{1}{Comprendre un exposant}\demarche{Relier l’écriture aux produits et aux inverses.}
Le premier produit vaut $\res{7^4}$. L’exposant négatif donne un inverse : $2^{-3}=1/2^3=\res{1/8}$. La base 9 est non nulle, donc $\res{9^0=1}$.
\end{corrige}\end{blocLivre}
\begin{blocLivre}{cor-2-2}\begin{corrige}{2}{Parenthèses et signe}\demarche{Déterminer la base exacte avant de calculer.}
Le carré de $-5$ vaut $\res{25}$, tandis que $-5^2=-(5^2)=\res{-25}$. Une puissance impaire de $-2$ est négative : $\res{-8}$ ; une puissance paire est positive : $\res{16}$.
\end{corrige}\end{blocLivre}
\begin{blocLivre}{cor-2-3}\begin{corrige}{3}{Règles de puissances}\demarche{Identifier la règle correspondant à chaque opération.}
Pour un produit de même base, on additionne les exposants : $\res{3^6}$. Pour le quotient, on les soustrait : $\res{5^4}$. Deux puissances de même exposant donnent $(2\times5)^3=\res{10^3}$. Enfin $10^{-2-4}=\res{10^{-6}}$.
\end{corrige}\end{blocLivre}
\begin{blocLivre}{cor-2-4}\begin{corrige}{4}{Écriture scientifique}\demarche{Contrôler le coefficient et le nombre de déplacements.}
On déplace la virgule jusqu’à obtenir un coefficient entre 1 et 10 : $\res{7{,}4\times10^7}$ et $\res{9{,}2\times10^{-5}}$. Le signe négatif de l’exposant traduit ici un nombre inférieur à 1.
\end{corrige}\end{blocLivre}
\begin{blocLivre}{cor-2-5}\begin{corrige}{5}{Quelques racines}\demarche{Utiliser les carrés connus en conservant le signe positif.}
On cherche des nombres positifs dont les carrés sont les données : $\res{12}$, $\res{0{,}9}$ et $\res{7/8}$. Par exemple $(7/8)^2=49/64$.
\end{corrige}\end{blocLivre}
\begin{blocLivre}{cor-2-6}\begin{corrige}{6}{Produit scientifique}\demarche{Regrouper puis normaliser le coefficient.}
Les coefficients donnent 24 et les puissances $10^{8-5}=10^3$. On obtient $24\times10^3=\res{2{,}4\times10^4}$. Le coefficient 24 n’était pas encore compris entre 1 et 10.
\end{corrige}\end{blocLivre}
\begin{blocLivre}{cor-2-7}\begin{corrige}{7}{Quotient scientifique}\demarche{Soustraire l’exposant du dénominateur, même s’il est positif.}
Le quotient des coefficients vaut 4 et celui des puissances vaut $10^{-3-5}=10^{-8}$. Le résultat est $\res{4\times10^{-8}}$. On vérifie qu’il est petit, puisque l’on divise un petit nombre par un grand.
\end{corrige}\end{blocLivre}
\begin{blocLivre}{cor-2-8}\begin{corrige}{8}{Des préfixes}\demarche{Remplacer le préfixe par sa puissance de 10, puis normaliser.}
Giga vaut $10^9$ : $\res{4{,}7\times10^9}$ octets. Nano vaut $10^{-9}$ : $250\times10^{-9}=\res{2{,}5\times10^{-7}\text{ m}}$. Micro vaut $10^{-6}$ : $0{,}6\times10^{-6}=\res{6\times10^{-7}\text{ s}}$. Les coefficients sont ensuite ramenés entre 1 et 10.
\end{corrige}\end{blocLivre}
\begin{blocLivre}{cor-2-9}\begin{corrige}{9}{Une équation $x^2=a$}\demarche{Examiner d’abord le signe du second membre.}
Les ensembles de solutions sont $\res{\{-7\,;7\}}$, $\res{\{-\sqrt7\,;\sqrt7\}}$, $\res{\{0\}}$ et $\res{\varnothing}$. Un carré n’est jamais négatif, ce qui explique le dernier cas.
\end{corrige}\end{blocLivre}
\begin{blocLivre}{cor-2-10}\begin{corrige}{10}{Réfuter une règle}\demarche{Choisir un contre-exemple dont les deux calculs sont exacts.}
Prendre $a=9$ et $b=16$ : $\sqrt{9+16}=\sqrt{25}=5$, tandis que $\sqrt9+\sqrt{16}=3+4=7$. Comme $5\ne7$, cet exemple suffit à réfuter l’affirmation générale.
\end{corrige}\end{blocLivre}
\begin{blocLivre}{cor-2-11}\begin{corrige}{11}{Une population qui double}\demarche{Compter les périodes avant d’appliquer la puissance.}
Deux heures représentent 120 minutes, donc six doublements. Le nombre prévu est $800\times2^6=\res{51200}$, soit $\res{5{,}12\times10^4}$. On multiplie six fois par 2 ; on n’ajoute pas six fois 800.
\end{corrige}\end{blocLivre}
\begin{blocLivre}{cor-2-12}\begin{corrige}{12}{Une masse minuscule}\demarche{Diviser puis convertir sans perdre les puissances de 10.}
On divise la masse totale par le nombre de particules : $1{,}8/(4\times10^{12})=0{,}45\times10^{-12}=\res{4{,}5\times10^{-13}\text{ g}}$. Comme $1\text{ g}=10^{-3}\text{ kg}$, cela vaut $\res{4{,}5\times10^{-16}\text{ kg}}$.
\end{corrige}\end{blocLivre}
\begin{blocLivre}{cor-2-13}\begin{corrige}{13}{Une distance astronomique}\demarche{Utiliser des unités compatibles avant de multiplier.}
La durée vaut $8\times60+20=500$ s. La distance est $3\times10^5\times500=\res{1{,}5\times10^8\text{ km}}$. On multiplie par $10^3$ pour convertir en mètres : $\res{1{,}5\times10^{11}\text{ m}}$.
\end{corrige}\end{blocLivre}
\begin{blocLivre}{cor-2-14}\begin{corrige}{14}{Choisir le côté d’une dalle}\demarche{Distinguer la longueur exacte et la mesure arrondie.}
Le côté $c$ vérifie $c^2=2$ avec $c>0$, donc $\res{c=\sqrt2\text{ m}}$, soit environ 1,414213 m. Au millimètre, on obtient $\res{1{,}414\text{ m}}$. Remplacer par 1,4 m donnerait $1{,}4^2=1{,}96$ m$^2$, différent de 2 : l’approximation ne conserve pas l’aire exacte.
\end{corrige}\end{blocLivre}
\begin{blocLivre}{cor-2-15}\begin{corrige}{15}{Un cercle et trois droites}\demarche{Se ramener à une équation $x^2=a$ et discuter selon le signe de $a$.}
Pour $y=3$ : $x^2=25-9=16$, donc deux points d’abscisses $\res{-4\text{ et }4}$. Pour $y=5$ : $x^2=0$, un seul point, $\res{(0\,;5)}$ : la droite touche le cercle. Pour $y=6$ : $x^2=-11$, $\res{\text{aucun point}}$. Le nombre d’intersections dépend de la distance de la droite au centre, comparée au rayon.
\label{fig-cor-2-15}\figureGeo[corrige]{cercle-droites}
\end{corrige}\end{blocLivre}
